% Coronary segments: the choice of arteries, their division into proximal, mid and distal
% segments, and a check of the division. Short main text (user, 2026-11-02): the exact division,
% thresholds and the landmark table are in appendix/B_segment_split_details.tex (sec:app-split).
% Placed after 02_population_outcome. Humanised 2026-11-02 in the style of clinical_background_v2.
% Does not repeat: tracer labels only the artery, Medis by segment (04a, sec:method-medis);
% what the visual analysis and the Medis trees and contours are (data/02, data/03).
% Division: CGPS_data/split_trunks.py (landmark_cuts). Check: CGPS_data/thesis_final/split_check.py
% --trees tracer, section 2 (contours by location vs the visual BL grade: a contour whose majority
% segment is not graded plaque = miss); reference run --trees medis.
% Red clarification notes; needs \usepackage{xcolor} in the preamble. Remove before submission.
% 2026-11-09: "further down, in the LAD and LCx" -> "at the start of the LAD and LCx, just beyond the
%   bifurcation", to agree with sec:method-features (angle and ratios enter only the proximal LAD
%   and LCx). CHECK: the review now ties the LM angle to low WSS at the bifurcation (Garcha,
%   simulations) and plaque to bifurcation outer walls (van der Giessen), but no study in it relates
%   the LM angle itself to disease; the study behind the original sentence (cui2017bifurcation?) was
%   cut from the review. Soften to "low WSS" or restore the source.
% 2026-11-20 (user: "move segment_split to there; now the segment split only affects curvature"): moved
%   from 02b_segment_split.tex (own section "Analysed Arteries and Segments", after the cohort) into the
%   features section, before sec:method-curvature. The angle and asymmetry use the LAD and LCx paths from
%   O (split_trunks.split), not the segment boundaries (split_trunks.landmark_cuts), and their outcome
%   (proximal LAD or LCx) comes from the visual reads, so only curvature depends on the division. Section
%   heading removed; Choice of Arteries (sec:method-population-arteries) moved to 02_population_outcome;
%   Division -> subsection, Validation -> subsubsection; "the features" / "features" -> curvature in
%   the first and last sentence. Labels unchanged.
\providecommand{\clarify}[1]{\textcolor{red}{\textbf{Clarification needed:} #1}}
\subsection{Division into Segments}
\label{sec:method-split}

Geometry is theorised to influence plaque formation through low wall shear stress, which is a local force. 
A curvature averaged over a whole artery would risk blurring these local differences, so curvature was computed at the same fine resolution as the outcome.
Each of the three arteries was therefore divided into a proximal, a mid and a distal segment, giving nine segments per participant (Figure~\ref{fig:method-segments}).

The visual analysis does not record where one segment ends and the next begins, so a landmark division method was developed and used in this thesis to place the boundaries.
Where no such branch was available, the boundary was placed at a fixed fraction of the distance along the artery.
The full division is given in Appendix~\ref{sec:app-split}.

% Figure from segment_selection_schematic.py (made-up tree, no patient data).

\subsubsection{Validation of the Division}
\label{sec:method-split-validation}

To give an indication of the reliability of the division, the Medis lumen contours were projected onto the nearest points of the divided tracer artery, and the segment holding the largest part of each contour was looked up in the visual analysis.
If that segment was not graded with plaque at baseline, a miss was recorded.
Only the position of the contour was used, not the segment name Medis gave it, so the check does not depend on how Medis divided the artery.
The miss rate is reported in Section~\ref{sec:results-cohort}.
A slightly displaced boundary moves only a few points at either end of a segment, so curvature computed per segment was preferred, despite this uncertainty, to curvature averaged over the whole artery.
